% SPDX-License-Identifier: GPL-3.0-or-later
% Arabic -- what the reading editions' preamble leaves to the language.
% Input by lib/frank-preamble.tex as \input{\FrankLib/lang/\BookLang.tex};
% docs/languages.md section 6 lists the hooks every file here must define.
% Copied from fa.tex: the same script, the same direction, one pass fewer.

% ---- the babel language and the switches ----------------------------------
\newcommand{\FrankLangName}{arabic}
\FrankRTLtrue            % chunk right, gloss left; \pw isolates with \babelsublr
\FrankHasBaretrue        % pass 3 is the same text without the tashkil
\FrankHasAltfalse        % no alternate face, so no fourth pass
\FrankHasReadingfalse
\FrankHasVertfalse

% ---- fonts ------------------------------------------------------------------
% From the registry's "tex" record: Noto Naskh Arabic travels with the toolbox
% (lib/fonts), Amiri and Vazirmatn are the fallbacks.  No alternate face:
% \FrankAltFont stays undefined, which is what \ifFrankHasAlt false promises.
\babelprovide[import=ar,onchar=ids fonts]{arabic}
\FrankPickFont{\FrankMainFontName}{ar}{main}{main_fallbacks}
\babelfont[arabic]{rm}[Script=Arabic,Language=Arabic,Renderer=HarfBuzz]{\FrankMainFontName}

% ---- the vocabulary line -----------------------------------------------------
% \vb{perfect}{sound (form)}{imperfect}{sound}{masdar}{sound}{meaning}
% The perfect is the dictionary's form, third person masculine singular, and
% its sound slot carries the verb's FORM after the transliteration, in
% brackets: kataba (I), ʿallama (II), arāda (IV), tarjama (Iq).  The number
% is what turns the other two slots into a pattern for every form but the
% first, so it rides where docs/lang/ar.md always put it and needs no slot of
% its own.  The imperfect is the third person masculine singular indicative,
% final -u and all: its vowel is the one thing a form-I verb will not tell
% you (yaktubu, yaqūlu, yaṣilu, yarā).  The masdar is the verbal noun that
% goes with the sense in the text -- وُصُول for "to arrive", not صِلَة -- and
% never the perfect written a second time: the first fixture edition did
% exactly that nine times over and printed "masdar رَأَى raʾā", and nothing
% caught it, because check_batch checks Arabic for fidelity and not for this.
% A verb with no masdar in use leaves the third pair blank, {}{}, and the
% core then prints neither the form nor the label (lib/frank-preamble.tex).
% The one extra, in a parenthesis after the meaning, is the preposition the
% verb governs where it is not a plain object:
%   \vb{وَصَلَ}{waṣala (I)}{يَصِلُ}{yaṣilu}{وُصُول}{wuṣūl}{to arrive (+ \pw{إِلَى})}
% Two alternatives share one + with a slash, (+ \pw{مِنْ}/\pw{عَنْ}); a second
% object has a second +, (+ \pw{لِ} + \pw{بِ}).  The gloss editor proposes the
% whole \vb from the dictionary (lib/verbs/ar.py): Wiktionary's head line gives
% the perfect and its form, the imperfect and the masdars, and the sense's
% +obj template the preposition -- but the masdar it proposes is the FIRST one
% listed, and for "to arrive" that is صِلَة: a draft, which the annotator
% corrects to the one above.
% The reader takes the same pair of labels from the registry
% (lib/languages.json, "vb_labels": ["impf.", "masdar"], and the three forms in
% words as "vb_forms"); the two must agree, or the PDF and the reader of one
% book label the same form differently.  lib/newlang.py --check compares them.
\newcommand{\FrankVbPres}{impf.}
\newcommand{\FrankVbPast}{masdar}

% ---- the two Lua filters ---------------------------------------------------
% frank_strip drops the tashkil the vowelled pass carries: the harakat range
% U+064B..U+0652 as for Persian, plus U+0670, the superscript alef of words
% like رَحْمٰن, which Persian text never writes and Arabic school editions do
% (the registry's `strip` for ar).  frank_latin maps the Arabic-Indic digits
% U+0660..U+0669 to ASCII for the destination names.  No #, % or backslash in
% a \directlua body (NOTES section 9 of the Persian book).
\directlua{
  function frank_strip(s)
    local t, n = {}, 0
    for _, c in utf8.codes(s) do
      if (c < 0x064B or c > 0x0652) and not (c == 0x0670) then n = n + 1 ; t[n] = utf8.char(c) end
    end
    return table.concat(t)
  end
  function frank_latin(s)
    local t, n = {}, 0
    for _, c in utf8.codes(s) do
      if c >= 0x0660 and c <= 0x0669 then c = c - 0x0660 + 0x30 end
      n = n + 1 ; t[n] = utf8.char(c)
    end
    return table.concat(t)
  end
}

% ---- the "How to read this" page -----------------------------------------------
\newcommand{\FrankHowTo}{%
Every passage comes three times.

\smallskip
\textbf{First} the whole sentence with the tashkil written in — fatha, kasra,
damma, sukun, shadda and the tanwin, the way a school edition or the Qur'an is
pointed — and nothing else. Try to read it. Get as far as you can and do not
worry about what you cannot get; the point is the attempt, made before any
help arrives.

\smallskip
\textbf{Second} the same sentence in small chunks, Arabic on the right, and on
the left three lines: how it sounds, then how to look each word up, then what
it means. Now you find out how much you had right.

\smallskip
\textbf{Third} the same passage again without the vowel marks, the way Arabic
is really written and printed. You have read it twice, so it should carry
you; the article, the long vowels and the shape of the word are what you are
learning to read by.

\smallskip
Verbs are given as the perfect (the dictionary form, third person masculine
singular), then the imperfect — the vowel of the imperfect is the one you
cannot guess — then the verbal noun (masdar), then the meaning. The Roman
numeral after the first transliteration is the verb's form: \textit{I} is the
plain verb, \textit{II} doubles its middle letter (\textit{ʿallama}, to
teach), \textit{X} begins \textit{ista-}, and every form but the first makes
its imperfect, and nearly always its masdar, to a pattern — so it is only a
verb of form \textit{I} whose three words have to be learnt one by one. A
preposition in brackets after the meaning is the one the verb takes:
\textit{to arrive} (+ \pw{إِلَى}) is to arrive \emph{at} somewhere; two with a
slash between them are two ways of saying it, and a second + is a second
object. Nouns that are broken plurals are given with their singular.

\smallskip
\textit{\=a}, \textit{\=\i}, \textit{\=u} are the long vowels; \textit{ʾ} is
the hamza and \textit{ʿ} the \textit{ʿayn}; a dot under a letter
(\textit{ṣ ḍ ṭ ẓ ḥ}) marks the emphatic and pharyngeal sounds, and
\textit{ṯ ḏ ḫ ġ š} are one letter each (\textit{th, dh, kh, gh, sh} in other
books); a doubled consonant is the shadda. The article is written
\textit{al-} with a hyphen and as it is heard: \textit{aš-šams},
\textit{an-nūr}, but \textit{al-qamar}.
}
